\documentclass[11pt,a4paper]{article}
\usepackage{fontspec,amsmath,amssymb,amsthm,microtype}
\setmainfont{lmroman10-regular.otf}[BoldFont=lmroman10-bold.otf,ItalicFont=lmroman10-italic.otf,BoldItalicFont=lmroman10-bolditalic.otf]
\usepackage{xcolor}
\usepackage[most]{tcolorbox}
\usepackage{mathrsfs}
\providecommand{\lt}{<}
\providecommand{\gt}{>}
\setmonofont{DejaVuSansMono.ttf}[BoldFont=DejaVuSansMono-Bold.ttf,ItalicFont=DejaVuSansMono-Oblique.ttf,BoldItalicFont=DejaVuSansMono-BoldOblique.ttf,Scale=0.86]
\newfontfamily\leanfallback{STIXTwoMath-Regular.otf}[Scale=0.86]
\definecolor{leanrule}{gray}{0.62}
\definecolor{leanback}{gray}{0.962}
\definecolor{leantitle}{gray}{0.22}
\definecolor{leanmuted}{gray}{0.38}
\newtcolorbox{leanbox}[1]{enhanced,breakable,colback=leanback,colframe=leanrule,
  boxrule=0.4pt,arc=0.6mm,left=2.4mm,right=2.4mm,top=1.8mm,bottom=1.8mm,
  before skip=0.8em,after skip=1em,fonttitle=\small\bfseries,coltitle=leantitle,
  colbacktitle=leanback,titlerule=0.3pt,title={#1}}
\newlength{\leancharwidth}
\newenvironment{leancode}{\par\smallskip\begingroup\ttfamily\footnotesize\raggedright
  \settowidth{\leancharwidth}{M}\setlength{\parskip}{0pt}\setlength{\parindent}{0pt}
  \frenchspacing}{\par\endgroup\smallskip}
\newcommand{\leanline}[2]{\hangindent=\dimexpr#1\leancharwidth+4\leancharwidth\relax
  \hangafter=1\noindent\hspace*{#1\leancharwidth}#2\par}
\newcommand{\leansource}[1]{\par{\footnotesize\color{leanmuted}#1}\par}
\newcommand{\leanlabel}[1]{\par\smallskip\noindent{\small\bfseries #1}\par}
\usepackage[margin=25mm]{geometry}
\usepackage[unicode,hidelinks]{hyperref}
\hypersetup{pdfauthor={LangevinLean Project}}
\setlength{\emergencystretch}{3em}
\setlength{\parskip}{0.4em}
\setlength{\parindent}{0pt}
\allowdisplaybreaks[2]
\raggedbottom
\title{\texorpdfstring{Integral identities and stationary laws of the second-order reference segments}{Integral identities and stationary laws of the second-order reference segments}}
\author{LangevinLean Project}
\date{}
\begin{document}\maketitle
Draft mathematical exposition; these notes may contain mistakes.


Each local second-order comparison uses a selected kinetic solution on a specified Brownian interval. This section exposes its continuous integral equation, exact endpoint stitching, stationary phase law, force-mild formula, and kinematic identity.



The global stationary reference supplies every local path equation, endpoint identity, and moment marginal needed by the corrected-step analysis. No separate reference-law premise remains at each numerical step.


\section{\texorpdfstring{Obtain continuity and endpoint matching from the solution constructor}{Obtain continuity and endpoint matching from the solution constructor}}


The reference segment \(Z_n(s)\), \(0\le s\le h\), is the selected continuous solution starting at the preceding reference state and driven by the Brownian increment beginning at \(nh\). The constructor proves continuity, \(Z_n(0)=R_n\), \(Z_n(h)=R_{n+1}\), and the kinetic integral equation \[Z_n(t)-Z_n(0)=-\int_0^tb(Z_n(s))\,ds+(\sqrt{2\gamma}\,\Delta W_n(t),0).\] The proof specializes these already constructed solutions with the lifted gradient and unit inverse mass. No local path or endpoint-matching assumption is introduced.



The selected integral flow gives continuity, the starting phase, and endpoint matching of every reference and auxiliary segment. Restart uniqueness links them to the global path where appropriate. These facts supply the structural hypotheses of the second-order local analysis.

\begin{leanbox}{Lean: \texttt{referenceSegment\_\allowbreak{}spec}}
\leansource{\texttt{referenceSegment\_\allowbreak{}spec} \quad KineticLangevin/\allowbreak{}CanonicalSecondOrderSegments.lean, lines 12–24}
\begin{leancode}
\leanline{0}{theorem\ referenceSegment\_spec\ (c\ :\ ULA.Coupling\ μ\ π)}
\leanline{4}{(hπ\ :\ MemLp\ (id\ :\ H\ ι\ →\ H\ ι)\ 2\ π)\ (f\ :\ H\ ι→ℝ)\ (M\ γ\ :\ ℝ≥0)}
\leanline{4}{(hLip\ :\ LipschitzWith\ M\ (gradient\ f))\ (h\ :\ ℝ≥0)\ (n\ :\ ℕ)\ :}
\leanline{4}{Continuous\ (referenceSegment\ c\ hπ\ f\ M\ γ\ hLip\ h\ n)\ ∧}
\leanline{6}{referenceSegment\ c\ hπ\ f\ M\ γ\ hLip\ h\ n\ 0=referenceChain\ c\ hπ\ f\ M\ γ\ hLip\ h\ n\ ∧}
\leanline{6}{referenceSegment\ c\ hπ\ f\ M\ γ\ hLip\ h\ n\ h=referenceChain\ c\ hπ\ f\ M\ γ\ hLip\ h\ (n+1)\ ∧}
\leanline{6}{∀\ t∈Icc\ 0\ (h\ :\ ℝ),referenceSegment\ c\ hπ\ f\ M\ γ\ hLip\ h\ n\ t-referenceSegment\ c\ hπ\ f\ M\ γ\ hLip\ h\ n\ 0=}
\leanline{8}{-(∫\ s\ in\ 0..t,kineticNegativeDrift\ (ULA.lpDrift\ (probability\ c)\ (gradient\ f)\ hLip)\ γ\ 1}
\leanline{10}{(referenceSegment\ c\ hπ\ f\ M\ γ\ hLip\ h\ n\ s))+(input\ c\ γ\ h\ n\ t,0)}
\end{leancode}
\leanlabel{How the Lean proof follows the argument}
\texttt{reference\allowbreak{}Segment\_\allowbreak{}spec} gives the formal counterpart of this step. Its remaining equalities are discharged by the simplifications displayed in the proof.
\end{leanbox}


Earlier results used here: \cite[Section~3.13.2]{LangevinLean2026ULASC}; \cite[Section~3.13.5]{LangevinLean2026ULASC}; \cite[Section~3.17.1]{LangevinLean2026ULASC}.

\section{\texorpdfstring{The mild and kinematic equations with the same fresh noise}{The mild and kinematic equations with the same fresh noise}}


Projecting the integral equation gives the velocity equation. The proved filtered specified Brownian response solves \(N(t)=\sqrt{2\gamma}\Delta W_n(t)-\gamma\int_0^tN(s)\,ds\); applying the integrating-factor theorem yields \[V_n(t)=R_0(t)-\int_0^te^{-\gamma(t-s)}\nabla f(X_n(s))\,ds.\] Here \(R_0\) is exactly the free-OU response used by the corrected numerical path. The positional projection gives \(X_n(t)-X_n(0)=\int_0^tV_n(s)\,ds\). The proof restricts the segment equation to each shorter duration, substitutes the exact filtered-input identity, simplifies unit inverse mass, and reorders the resulting additive terms.



The mild velocity and integrated-position identities follow from the local kinetic equation and filtered input. The same interval source generates all four correction responses. Thus the analytic corrected path and numerical map are compared under their common forcing.

\begin{leanbox}{Lean: \texttt{referenceSegment\_\allowbreak{}mild}, \texttt{referenceSegment\_\allowbreak{}kinematic}}
\leansource{\texttt{referenceSegment\_\allowbreak{}mild} \quad KineticLangevin/\allowbreak{}CanonicalSecondOrderSegments.lean, lines 27–44}
\begin{leancode}
\leanline{0}{theorem\ referenceSegment\_mild\ (c\ :\ ULA.Coupling\ μ\ π)}
\leanline{4}{(hπ\ :\ MemLp\ (id\ :\ H\ ι\ →\ H\ ι)\ 2\ π)\ (f\ :\ H\ ι→ℝ)\ (M\ γ\ :\ ℝ≥0)}
\leanline{4}{(hLip\ :\ LipschitzWith\ M\ (gradient\ f))\ (h\ :\ ℝ≥0)\ (n\ :\ ℕ)\ (t\ :\ ℝ)\ (ht\ :\ t∈Icc\ 0\ (h\ :\ ℝ))\ :}
\leanline{4}{(referenceSegment\ c\ hπ\ f\ M\ γ\ hLip\ h\ n\ t).1=}
\leanline{6}{freeOUResponse\ (probability\ c)\ (inputLift\ c\ γ)}
\leanline{8}{(referenceSegment\ c\ hπ\ f\ M\ γ\ hLip\ h\ n\ 0).1\ γ\ ((n\ :\ ℝ≥0)*h)\ 0\ t-}
\leanline{6}{dampedIntegral\ γ\ (fun\ s\ ↦\ ULA.lpDrift\ (probability\ c)\ (gradient\ f)\ hLip}
\leanline{8}{(referenceSegment\ c\ hπ\ f\ M\ γ\ hLip\ h\ n\ s).2)\ t}
\end{leancode}
\leansource{\texttt{referenceSegment\_\allowbreak{}kinematic} \quad KineticLangevin/\allowbreak{}CanonicalSecondOrderSegments.lean, lines 47–55}
\begin{leancode}
\leanline{0}{theorem\ referenceSegment\_kinematic\ (c\ :\ ULA.Coupling\ μ\ π)}
\leanline{4}{(hπ\ :\ MemLp\ (id\ :\ H\ ι\ →\ H\ ι)\ 2\ π)\ (f\ :\ H\ ι→ℝ)\ (M\ γ\ :\ ℝ≥0)}
\leanline{4}{(hLip\ :\ LipschitzWith\ M\ (gradient\ f))\ (h\ :\ ℝ≥0)\ (n\ :\ ℕ)\ (t\ :\ ℝ)\ (ht\ :\ t∈Icc\ 0\ (h\ :\ ℝ))\ :}
\leanline{4}{(referenceSegment\ c\ hπ\ f\ M\ γ\ hLip\ h\ n\ t).2-(referenceSegment\ c\ hπ\ f\ M\ γ\ hLip\ h\ n\ 0).2=}
\leanline{6}{∫\ s\ in\ 0..t,(referenceSegment\ c\ hπ\ f\ M\ γ\ hLip\ h\ n\ s).1}
\end{leancode}
\leanlabel{How the Lean proof follows the argument}
\texttt{reference\allowbreak{}Segment\_\allowbreak{}mild} gives the formal counterpart of this step. The main intermediate facts are \texttt{hs}, \texttt{hm} in that order. Its remaining equalities are discharged by the simplifications displayed in the proof.




\texttt{reference\allowbreak{}Segment\_\allowbreak{}kinematic} gives the formal counterpart of this step. The main intermediate facts are \texttt{hs} in that order. Named proof applications include \texttt{kinetic\_\allowbreak{}integral\_\allowbreak{}kinematic}; the adjacent proof shows their hypotheses and order of use.
\end{leanbox}


Earlier results used here: \cite[Section~3.13.2]{LangevinLean2026ULASC}; \cite[Section~3.13.5]{LangevinLean2026ULASC}; \cite[Section~3.17.1]{LangevinLean2026ULASC}.

\section{\texorpdfstring{Transfer global stationarity to every local segment point}{Transfer global stationarity to every local segment point}}


The reference segment is already proved equal to the single global reference solution at time \(nh+t\). Choose horizon \(T=(n+1)h\); the segment time lies in \([0,T]\). The argument rewrites the local solution using this stitching identity and applies the proved unconditional Gibbs stationarity of the global constructed solution. Thus the full phase law at every segment point, including both endpoints, is \(N(0,I)\otimes\pi\). This supplies the stationary moment hypotheses for the local second-order Taylor estimate.



Global phase stationarity gives the target position and Gaussian velocity marginals at every local time. Exact segment identification transfers those laws without a new stationarity assumption at each restart. The required moment bounds therefore hold throughout the local integration interval.

\begin{leanbox}{Lean: \texttt{referenceSegment\_\allowbreak{}stationary}}
\leansource{\texttt{referenceSegment\_\allowbreak{}stationary} \quad KineticLangevin/\allowbreak{}CanonicalSecondOrderSegments.lean, lines 58–75}
\begin{leancode}
\leanline{0}{theorem\ referenceSegment\_stationary}
\leanline{4}{(f\ :\ H\ ι→ℝ)\ (m\ :\ ℝ)\ (hm\ :\ 0<m)\ (M\ γ\ :\ ℝ≥0)}
\leanline{4}{(hf\ :\ Differentiable\ ℝ\ f)\ (hsc\ :\ StrongConvexOn\ univ\ m\ f)\ (hLip\ :\ LipschitzWith\ M\ (gradient\ f))}
\leanline{4}{(c\ :\ ULA.Coupling\ μ\ (ULA.gibbsMeasure\ f))}
\leanline{4}{(hπ\ :\ MemLp\ (id\ :\ H\ ι→H\ ι)\ 2\ (ULA.gibbsMeasure\ f))}
\leanline{4}{(h\ :\ ℝ≥0)\ (n\ :\ ℕ)\ (t\ :\ ℝ)\ (ht\ :\ t∈Icc\ 0\ (h\ :\ ℝ))\ :}
\leanline{4}{HasLaw\ (lpPhaseCoe\ (probability\ c)\ (referenceSegment\ c\ hπ\ f\ M\ γ\ hLip\ h\ n\ t))}
\leanline{6}{(kineticGibbsTarget\ f)\ (probability\ c)}
\end{leancode}
\leanlabel{How the Lean proof follows the argument}
\texttt{reference\allowbreak{}Segment\_\allowbreak{}stationary} gives the formal counterpart of this step. The main intermediate facts are \texttt{T}, \texttt{hn}, \texttt{h\allowbreak{}T}, \texttt{htime} in that order. Named proof applications include \texttt{global\allowbreak{}Reference\allowbreak{}Flow\_\allowbreak{}stationary}; the adjacent proof shows their hypotheses and order of use.
\end{leanbox}


Earlier results used here: \cite[Section~3.17.1]{LangevinLean2026ULASC}; \cite[Section~3.26.3]{LangevinLean2026ULASC}.

\begin{thebibliography}{ULA-LSI}
\bibitem[ULA-SC]{LangevinLean2026ULASC}
LangevinLean Project. \emph{Unadjusted Langevin under Strong Convexity}. LL-01, 2026. Draft report, content revision d34ead06987c. \url{https://odakyildiz.github.io/LangevinLean/docs/pdf/UnadjustedLangevin/FullProof.pdf}.
\end{thebibliography}

\end{document}
